\section{FAC EVIDENCE: RECURRENCE, NEW FINDINGS, AND MIXED CASES}\label{sec:fac}

Prior work has established that Single Audit findings recur and that a
reported finding depends jointly on the condition existing and on the auditor
detecting it
\citep{keating2005,petrovits2011,waymire2018,gao2002,gao2024,waddell2023}.
What is added is the decomposition of what follows a prior finding into
recurrence, newly recorded findings and mixed cases.

\subsection*{Data and sample}

The data are the public Federal Audit Clearinghouse bulk files; the unit of
analysis is the auditee--program--year, and the window is fiscal years 2016
through 2023, inside a single audit-threshold regime \citep{omb2024}. Of
\nPanelRows{} program-years, \nResolvedRows{} have a resolved prior-year
observation of the same program, and Sample~A restricts further to the
\nSampleA{} in which the program was audited as a major program in both
years: \nNPrior{} with a prior-year finding and \nNNoPrior{} without.
Intervals are cluster bootstrap over \nNAuditees{} auditees.
Appendix~\ref{app:emp} gives the construction and the sample flow.

One construction decision follows from the paper's own argument. In
\nSthreeN{} program-years, \nSthreePct{} percent of the panel, the auditee
filed in the prior year but this program does not appear in that filing. They
are excluded rather than coded as ``no prior finding'', because the absence of
a record is not a clean result: they go on to record a finding
\nSthreePctFinding{} percent of the time, against \nStwoPctFinding{} percent
after a resolved clean prior year. ``Repeat'' is the auditor's own coding; a
validated historical link is the secondary definition
(Appendix~\ref{app:emp}).

\subsection*{What follows a prior finding}

Even among programs audited as major in both consecutive years, prior recorded
findings predict two distinct parts of the future finding record: explicit
recurrence, and additional newly recorded findings. Table~\ref{tab:fourstate}
is the main result. For context, a program-year following a clean prior year
records a finding \mPAnyNoPrior{} of the time and one following a prior-year
finding records one \mPAnyPrior{} of the time, a difference of \mDAny{}
$[\mDAnyLo{},\,\mDAnyHi{}]$ and a ratio of \nRRAny{}. The content is in
how that difference is made up.

A program-year lands in exactly one of four states: no finding ($0$); a repeat
finding and nothing else ($R$); a new finding and nothing else ($N$); or both
a repeat and a new finding ($B$). Write $\Delta$ for the difference in the
probability of a state between program-years with and without a prior
finding. Because the four states do not overlap and cover every possibility,
the differences add:
%
\begin{equation}\label{eq:four}
  \Delta_{\text{any finding}}
  \;=\;
  \Delta_{R}
  \;+\;
  \Delta_{N}
  \;+\;
  \Delta_{B} .
\end{equation}
%
The categories can be added without deciding where mixed cases belong,
because each mixed case is counted once, in $B$. In Sample~A, $\Delta_R =
+\nDRonly{}$, $\Delta_N = +\nDNonly{}$ and $\Delta_B = +\nDBoth{}$, adding to
$+\nDAny{}$. Every interval in the table excludes zero.

Most of the gap is repeat-involving, as an auditor would expect. The row that
is easy to read past is the one about new findings. Programs whose \emph{only}
outcome was a new finding account for $+\nDNonly{}$, a ratio of \nRRNonly{},
which looks like almost nothing. But ``new only'' does not answer the
question a reader is asking, which is whether the program recorded \emph{any}
newly recorded problem. That quantity is
%
\[
  \text{any new} \;=\; N + B ,
\]
%
because a period that recorded both a repeat and something new did record
something new. On that margin the difference is $+\nDAnyNew{}$
$[\nDAnyNewLo{},\,\nDAnyNewHi{}]$, from \mPAnyNewNoPrior{} to
\mPAnyNewPrior{}, a ratio of \nRRAnyNew{}. The two ratios, \nRRNonly{} and
\nRRAnyNew{}, differ by a factor of \nRRfactor{} and come from the same table.
The difference between them is definitional, not evidence about follow-up:
``new only'' excludes every period that also recorded a repeat, and ``any
new'' includes them. Periods land in $B$ for more than one reason, and these
data do not separate them: thorough follow-up of a repeat can put a program
there that would otherwise have shown only a new finding, and so can a program
simply having several distinct deficiencies.

The cleanest version counts findings rather than flagging programs. Every
individual finding is classified once, as a repeat or as new, so the
difference in the expected number of findings splits with no allocation rule
at all:
%
\begin{equation}\label{eq:count}
  \underbrace{+\nEtotal{}}_{\text{total}}
  \;=\;
  \underbrace{+\nErepeat{}}_{\text{repeat}}
  \;+\;
  \underbrace{+\nEnew{}}_{\text{new}} .
\end{equation}
%
A program-year following a prior-year finding records about one more finding
than one following a clean year, of which roughly seven tenths of a finding
is a repeat and three tenths is new.

The record does not identify how much of these associations reflects true
persistence of the underlying condition, stable differences between auditees,
condition-specific follow-up, broader scrutiny of programs with a history,
remediation combined with heavier detection, or mixtures of these. Each is
consistent with Table~\ref{tab:fourstate}. The FAC results illustrate the
identification problem; they do not identify its mechanism, and nothing here
should be cited as showing that observation feedback is the dominant
explanation.

\subsection*{A recurrence share is a convention, not a measurement}

The obvious way to compress Table~\ref{tab:fourstate} into one number is to
ask what share of the prior-finding association is recurrence. That statistic
is not uniquely defined when a program-period can contain both repeat and new
findings, because every mixed period, \nMixedShare{} percent of all
prior-finding outcomes, has to be assigned to one side. Assign none of them to
recurrence and the share is \nZetaLo{} percent. Assign all of them and it is
\nZetaHi{} percent. The summary moves \nZetaSwing{} percentage points, and
nothing in the underlying observations changes (Figure~\ref{fig:zeta}). The
general lesson: when outcome classes overlap, report mutually exclusive states
or additive counts before reporting an attributable share, and if a share is
reported anyway, state the convention and show the figure under the opposite
one. The convention-free statement is equation~\eqref{eq:four}:
$\nDAny{} = \nDRonly{} + \nDNonly{} + \nDBoth{}$.

\subsection*{A comparison the data do not support}

A natural next question is whether prior findings predict problems in
requirement areas that never had one. The Clearinghouse records the
requirement category attached to every reported finding. It does not record
the universe of requirements that applied to, and were tested in, each
program-year; that universe is defined by the Compliance Supplement and is
external to these data. A ``first finding in a previously clean area'' rate
therefore has no defensible denominator, and without one the comparison goes
wrong because the two history groups do not have the same opportunity to
produce the outcome. Take two requirement areas, each with a 0.20 chance of a
finding this year, independent of anything before. A program with neither
area previously flagged has a $1 - 0.80^2 = 0.36$ chance of a finding in an
area new to it; a program with one area already flagged has only the other
available, so its chance is $0.20$. Nothing underneath differs. In a
constructed twelve-area case in which nothing depends on history, the naive
new-area risk ratio is \mNullCategoryRR{}: prior-finding programs look about
thirty percent \emph{less} likely to have a new kind of problem, from unequal
opportunity alone. No denominator was invented, and the comparison is not
made.

\subsection*{Robustness}

Table~\ref{tab:specs} in Appendix~\ref{app:emp} reports every specification.
Restricting to transitions with stable Clearinghouse provenance moves the
any-finding gap from \nDAny{} to \nProvDeltaAny{} and leaves any-new at
\nProvDeltaAnyNew{}. Substituting the validated historical link for the
auditor's repeat flag moves the repeat-only gap to \nVlinkDeltaRonly{} and
the both-kinds gap to \nVlinkDeltaBoth{}, leaves any-new unchanged, and opens
a fifth outcome state worth \nVlinkFifth{} of the gap, given its own column
there. A prespecified three-year history rule, under which ``no prior
finding'' means three observed clean years, leaves \nHtwoSampleA{}
program-years and attenuates the gaps to \nHtwoDeltaAny{} for any finding and
\nHtwoDeltaAnyNew{} for any new. Across all five specifications the any-new
margin moves least.

\begin{table}[tbp]
\centering
\caption{What follows a prior-year finding.}
\label{tab:fourstate}
\footnotesize
\setlength{\tabcolsep}{4.5pt}
\begin{tabular}{lccccr}
\toprule
& \multicolumn{2}{c}{Chance this year} & & & \\
\cmidrule(lr){2-3}
Outcome for the program-year & \shortstack{no prior\\finding} & \shortstack{prior\\finding}
& Difference & 95\% interval & Ratio \\
\midrule
No finding at all & 0.9234 & 0.4217 & $-0.5017$ & $[-0.5104,\, -0.4927]$ & 0.46 \\
A repeat finding only & 0.0015 & 0.3088 & $+0.3073$ & $[0.3008,\, 0.3140]$ & --- \\
A new finding only & 0.0744 & 0.1200 & $+0.0456$ & $[0.0416,\, 0.0498]$ & 1.61 \\
Both a repeat and a new finding & 0.0007 & 0.1495 & $+0.1488$ & $[0.1416,\, 0.1561]$ & --- \\
\midrule
Any finding & 0.0766 & 0.5783 & $+0.5017$ & $[0.4927,\, 0.5104]$ & 7.55 \\
Any repeat finding & 0.0022 & 0.4583 & $+0.4561$ & $[0.4458,\, 0.4656]$ & --- \\
\textbf{Any new finding} (new only $+$ both) & 0.0751 & 0.2696 & $+0.1945$ & $[0.1870,\, 0.2021]$ & 3.59 \\
\midrule
\multicolumn{6}{l}{\emph{Findings counted rather than programs flagged --- this needs no allocation rule:}} \\
Findings recorded, total & 0.1029 & 1.1712 & $+1.0683$ & $[1.0209,\, 1.1182]$ & \\
\quad of which repeat & 0.0027 & 0.7468 & $+0.7441$ & $[0.7084,\, 0.7807]$ & \\
\quad of which new & 0.1002 & 0.4245 & $+0.3243$ & $[0.3076,\, 0.3422]$ & \\
\bottomrule
\end{tabular}
\tablenotes{Sample A: 206,926 program-years
in which the same program was audited as a major program in both years, 39,067 with
a finding in the prior year and 167,859 without. The four outcomes in the first panel
are mutually exclusive and, under the auditor's repeat flag used here, exhaust
the possibilities, so their risk differences add to the difference in ``any
finding'' exactly. Intervals are 95 percent
cluster bootstrap, 1,000 replications resampling 45,952 auditees with their full
histories, seed 20260918. Every interval excludes zero. The row that is easy to miss
is ``any new finding'': counting only the periods whose \emph{sole} outcome was
new understates it by roughly a factor of 2.23 on the ratio scale, because ``new
only'' excludes every period that also recorded a repeat. A period can land in
``both'' because repeat follow-up is thorough, or because the program has
several distinct deficiencies; these data do not separate the two. The
ratio column is left blank for the three repeat-involving rows.
Repeat-involving baseline probabilities are very small in this sample, making
their ratios sensitive to small denominators; risk differences are reported
instead.}
\end{table}


\begin{figure}[tbp]
\centering
\includegraphics[width=0.78\textwidth]{ex7_zeta.pdf}
\caption{\capZeta}
\label{fig:zeta}
\end{figure}
